\documentclass[10pt,a4paper]{amsart}
\usepackage{amsmath,amssymb,amsthm,graphicx}
\usepackage{verbatim}
\newenvironment{reference}{\comment}{\endcomment}
%\newenvironment{reference}{}{}
\newenvironment{slogan}{\comment}{\endcomment}
\newenvironment{history}{\comment}{\endcomment}

% For commutative diagrams we use Xy-pic
\usepackage[all]{xy}

% We use 2cell for 2-commutative diagrams.
\xyoption{2cell}
% Theorem environments.
%
\theoremstyle{plain}
\newtheorem{theorem}[subsection]{Theorem}
\newtheorem{proposition}[subsection]{Proposition}
\newtheorem{lemma}[subsection]{Lemma}

\theoremstyle{definition}
\newtheorem{definition}[subsection]{Definition}
\newtheorem{example}[subsection]{Example}
\newtheorem{exercise}[subsection]{Exercise}
\newtheorem{situation}[subsection]{Situation}

\theoremstyle{remark}
\newtheorem{remark}[subsection]{Remark}
\newtheorem{remarks}[subsection]{Remarks}

\numberwithin{equation}{subsection}

% Macros
%
\def\lim{\mathop{\mathrm{lim}}\nolimits}
\def\colim{\mathop{\mathrm{colim}}\nolimits}
\def\Spec{\mathop{\mathrm{Spec}}}
\def\Hom{\mathop{\mathrm{Hom}}\nolimits}
\def\Ext{\mathop{\mathrm{Ext}}\nolimits}
\def\SheafHom{\mathop{\mathcal{H}\!\mathit{om}}\nolimits}
\def\SheafExt{\mathop{\mathcal{E}\!\mathit{xt}}\nolimits}
\def\Sch{\mathit{Sch}}
\def\Mor{\mathop{\mathrm{Mor}}\nolimits}
\def\Ob{\mathop{\mathrm{Ob}}\nolimits}
\def\Sh{\mathop{\mathit{Sh}}\nolimits}
\def\NL{\mathop{N\!L}\nolimits}
\def\CH{\mathop{\mathrm{CH}}\nolimits}
\def\proetale{{pro\text{-}\acute{e}tale}}
\def\etale{{\acute{e}tale}}
\def\QCoh{\mathit{QCoh}}
\def\Ker{\mathop{\mathrm{Ker}}}
\def\Im{\mathop{\mathrm{Im}}}
\def\Coker{\mathop{\mathrm{Coker}}}
\def\Coim{\mathop{\mathrm{Coim}}}

% Boxtimes
%
\DeclareMathSymbol{\boxtimes}{\mathbin}{AMSa}{"02}

%
% Macros for moduli stacks/spaces
%
\def\QCohstack{\mathcal{QC}\!\mathit{oh}}
\def\Cohstack{\mathcal{C}\!\mathit{oh}}
\def\Spacesstack{\mathcal{S}\!\mathit{paces}}
\def\Quotfunctor{\mathrm{Quot}}
\def\Hilbfunctor{\mathrm{Hilb}}
\def\Curvesstack{\mathcal{C}\!\mathit{urves}}
\def\Polarizedstack{\mathcal{P}\!\mathit{olarized}}
\def\Complexesstack{\mathcal{C}\!\mathit{omplexes}}
% \Pic is the operator that assigns to X its picard group, usage \Pic(X)
% \Picardstack_{X/B} denotes the Picard stack of X over B
% \Picardfunctor_{X/B} denotes the Picard functor of X over B
\def\Pic{\mathop{\mathrm{Pic}}\nolimits}
\def\Picardstack{\mathcal{P}\!\mathit{ic}}
\def\Picardfunctor{\mathrm{Pic}}
\def\Deformationcategory{\mathcal{D}\!\mathit{ef}}
\usepackage{cleveref}
\crefname{originallemma}{original lemma}{original lemmas}
\Crefname{originallemma}{Original lemma}{Original lemmas}
\makeatletter
\newlabel{lemma-identify-chow-for-smooth}{{0FC2}{source}{}{sourcepointer-1}{}}
\newlabel{lemma-identify-chow-for-smooth@cref}{{[originallemma][0][]0FC2}{[1][0][]source}{}{}{}}
\makeatother
\crefname{originallemma}{original lemma}{original lemmas}
\Crefname{originallemma}{Original lemma}{Original lemmas}
\makeatletter
\newlabel{lemma-lci-gysin-commutes}{{0FF6}{source}{}{sourcepointer-2}{}}
\newlabel{lemma-lci-gysin-commutes@cref}{{[originallemma][0][]0FF6}{[1][0][]source}{}{}{}}
\makeatother
\crefname{originallemma}{original lemma}{original lemmas}
\Crefname{originallemma}{Original lemma}{Original lemmas}
\makeatletter
\newlabel{lemma-lci-gysin-easy}{{0FF7}{source}{}{sourcepointer-3}{}}
\newlabel{lemma-lci-gysin-easy@cref}{{[originallemma][0][]0FF7}{[1][0][]source}{}{}{}}
\makeatother
\title{The intersection product with the fundamental class of an equidimensional regularly embedded subscheme is the Gysin pullback pushed forward}
\author{The Stacks Project Authors}
\date{stacks-project:chow.tex:local-snapshot-20261012; GNU FDL 1.2 or later}
\begin{document}
\maketitle
\noindent\textit{Editorial note.} All delivered mathematical text is the author's own, copied from the pinned Stacks Project chapter file chow.tex; no statement or proof step was rewritten, reordered or removed.\par\medskip
\noindent\textit{Editorial note.} Preamble provenance: the pinned chapter file opens with an input of the shared preamble and carries no preamble text of its own, so the project's single Stacks-wide preamble staged at sources/stacks-chow/preamble.tex is excerpted verbatim as two spans, lines 14-18 (Xy-pic setup) and lines 154-217 (the native theorem declarations and mathematical macros). Nothing was added, removed or redefined; the wrapper class and the amsmath/amssymb/amsthm/graphicx base packages are supplied by the assembler.\par\medskip
\noindent\textit{Editorial note.} The counted statement and proof contain no cross reference, no citation and no included file, so the delivered body is self contained.\par\medskip
\noindent\textit{Editorial note.} Canonical location of the counted proposition: Stacks tag 0FFE (sources/stacks-chow/pinned-original-source.tex line 14563).\par\medskip
\noindent\textit{Allowed standard prerequisites (editorial).}\begin{itemize}
\item Identification of Chow groups with operational Chow cohomology for a smooth scheme. Same chapter, Lemma labelled lemma-identify-chow-for-smooth (source line 14403), invoked at proof lines 14572-14573 and 14583. Supplies the representation alpha = c cap [X]\_n and the final translation back to the intersection product.
\item Commutation of the Gysin map with capping by a bivariant class. Same chapter, Lemma labelled lemma-lci-gysin-commutes (source line 13613), invoked at proof lines 14579-14580. Supplies i\textasciicircum{}!(c cap [X]\_n) = c cap i\textasciicircum{}![X]\_n, the second step of the chain.
\item The Gysin pullback of the fundamental class along a regular closed immersion. Same chapter, Lemma labelled lemma-lci-gysin-easy (source line 13643), invoked at proof line 14581. Supplies i\textasciicircum{}![X]\_d = [Y]\_e in CH\_*(Y); it is counted separately in this lane as S12-algebra-lci-gysin-easy.
\end{itemize}\medskip
\ifdefined\hypertarget\hypertarget{sourcepointer-1}{}\fi
\noindent\textit{Source reference.} Original lemma 0FC2: sources/stacks-chow/pinned-original-source.tex, source lines 14403--14403. Same-chapter pointer into the author environment lemma at source line 14403; invoked at proof line 14573 of the counted proof. It is retained only as a navigation pointer; every load-bearing step is inside the delivered span.\par\medskip
\ifdefined\hypertarget\hypertarget{sourcepointer-2}{}\fi
\noindent\textit{Source reference.} Original lemma 0FF6: sources/stacks-chow/pinned-original-source.tex, source lines 13613--13613. Same-chapter pointer into the author environment lemma at source line 13613; invoked at proof line 14580 of the counted proof. It is retained only as a navigation pointer; every load-bearing step is inside the delivered span.\par\medskip
\ifdefined\hypertarget\hypertarget{sourcepointer-3}{}\fi
\noindent\textit{Source reference.} Original lemma 0FF7: sources/stacks-chow/pinned-original-source.tex, source lines 13643--13643. Same-chapter pointer into the author environment lemma at source line 13643; invoked at proof line 14581 of the counted proof. It is retained only as a navigation pointer; every load-bearing step is inside the delivered span.\par\medskip
\begin{lemma}
\label{lemma-intersect-regularly-embedded}
Let $k$ be a field. Let $X$ be a scheme smooth over $k$. Let $i : Y \to X$ be
a regular closed immersion. Let $\alpha \in \CH_*(X)$. If $Y$ is
equidimensional of dimension $e$, then
$\alpha \cdot [Y]_e = i_*(i^!(\alpha))$ in $\CH_*(X)$.
\end{lemma}

\begin{proof}
After decomposing $X$ into connected components we may and do assume $X$
is equidimensional of dimension $d$. Write $\alpha = c \cap [X]_n$
with $x \in A^*(X)$, see Lemma \ref{lemma-identify-chow-for-smooth}. Then
$$
i_*(i^!(\alpha)) = i_*(i^!(c \cap [X]_n)) =
i_*(c \cap i^![X]_n) = i_*(c \cap [Y]_e) =
c \cap i_*[Y]_e = \alpha \cdot [Y]_e
$$
The first equality by choice of $c$. Then second equality by
Lemma \ref{lemma-lci-gysin-commutes}. The third because
$i^![X]_d = [Y]_e$ in $\CH_*(Y)$ (Lemma \ref{lemma-lci-gysin-easy}).
The fourth because bivariant classes commute with proper pushforward.
The last equality by Lemma \ref{lemma-identify-chow-for-smooth}.
\end{proof}

\end{document}
